% SPDX-License-Identifier: AGPL-3.0-or-later
% Commercial license available
% © Concepts 1996–2026 Miroslav Šotek. All rights reserved.
% © Code 2020–2026 Miroslav Šotek. All rights reserved.
% ORCID: 0009-0009-3560-0851
% Contact: www.anulum.li | protoscience@anulum.li
% SCPN-FUSION-CORE — runaway kinetic operator with an equivalent Rust kernel.
\documentclass[11pt]{article}

\usepackage[utf8]{inputenc}
\usepackage[T1]{fontenc}
\usepackage{amsmath}
\usepackage{booktabs}
\usepackage{microtype}
\usepackage{natbib}
\usepackage[colorlinks=true,linkcolor=blue,citecolor=blue,urlcolor=blue]{hyperref}
\usepackage[margin=1in]{geometry}
\bibliographystyle{unsrtnat}

\title{A Finite-Volume Runaway-Electron Kinetic Operator with a Bit-Identical
Rust Kernel}
\author{Miroslav \v{S}otek\\
\small ANULUM CH\&LI, Marbach SG, Switzerland, \url{https://www.anulum.li}\\
\small ORCID: \href{https://orcid.org/0009-0009-3560-0851}{0009-0009-3560-0851}\\
\small \href{mailto:protoscience@anulum.li}{protoscience@anulum.li}}
\date{Preprint --- 16 November 2026}

\begin{document}
\maketitle

\begin{abstract}
SCPN-Fusion-Core evolves a runaway-electron distribution on a grid of minor
radius, pitch and momentum with a conservative finite-volume operator.  The
operator returns each physical tendency separately: radial transport, electric
acceleration, collisional drag and diffusion, pitch scattering, cross
diffusion, synchrotron and bremsstrahlung losses, avalanche generation and an
external source.  It also evolves the total runaway density, with its own
budgets, alongside the distribution.  The operator has a NumPy implementation
and a Rust kernel behind the same Python interface.  On a parity case of
$8 \times 16 \times 32$ cells advanced by four steps of a three-stage
strong-stability-preserving Runge--Kutta scheme, the two backends return
bit-identical distributions, tendencies and runaway densities.  The integrated
moments differ by at most a relative $L_2$ of $5.4 \times 10^{-14}$.  On one
loaded workstation the Rust kernel ran 2.10 times faster than NumPy for this
case.  The coefficients of the parity case are synthetic and chosen to exercise
every term; it is not a physical runaway scenario.  We claim no equivalence
with the DREAM code and no portable speed-up.
\end{abstract}

\noindent\textbf{Keywords:} runaway electrons; kinetic equation; finite volume;
Rust; code equivalence; disruption

\section{Introduction}

Runaway electrons generated in tokamak disruptions are a central concern for
machine protection.  DREAM solves flux-surface-averaged transport equations
for the plasma density, temperature and poloidal flux, with a bounce-averaged
kinetic equation for the electrons, to simulate the cooling and the runaway
dynamics of a disruption self-consistently \citep{Hoppe2021}.  SCPN-Fusion-Core
contains a smaller kinetic operator for the runaway distribution, with a
NumPy implementation and a compiled Rust kernel.  This note establishes one
property of that operator: its two implementations are equivalent to rounding,
term by term.  A comparison with DREAM is a separate campaign and is not part
of this note.

\section{The operator}

The distribution $f$ lives on cells of minor radius $r$, pitch $\xi$ and
momentum $p$ (in units of $m_e c$).  The operator is written in conservative
finite-volume form: every advection and diffusion term is the divergence of
face fluxes, with upwind face values for advection.  It evaluates, and returns separately, the tendencies of radial
transport, electric acceleration (in momentum and in pitch), collisional drag
and momentum diffusion, pitch scattering, the two cross diffusions between
momentum and pitch, the synchrotron and bremsstrahlung losses, avalanche
generation and an external source.  Next to $f$ it evolves the total runaway
density per radial cell, with separate budgets for its radial transport,
avalanche generation and external source.

A solver advances the state with a three-stage strong-stability-preserving
Runge--Kutta scheme with a bounded step.  Every stage must stay finite, and
the runaway density must stay non-negative, or the solver stops.  The backend is chosen per call: the
same solver object runs either the NumPy operator or the Rust kernel through
its Python binding.

\section{Equivalence of the two backends}

The parity case uses 8 radial cells on $0 \le r \le 0.2$\,m, 16 pitch cells on
$-1 \le \xi \le 1$ and 32 momentum cells on $0.02 \le p \le 20$.  Every
coefficient is a constant multiplied by $(1 + i/n)$, with $i$ the flat index of
the face or cell and $n$ their number, so that each term varies in space and
none vanishes.  The run returns the state at 0, 50 and 100\,ns with a largest
step of 25\,ns, which is four steps.  These coefficients do not describe a
physical runaway scenario; they exist to exercise the code.

Table~\ref{tab:parity} lists the comparison.  The distribution, all nine
tendencies, their total, the runaway density and its avalanche and
external-source budgets are bit-identical.  Five quantities differ at rounding
level: the three moments of the distribution (density, current density and
kinetic-energy density), the radial-transport budget of the runaway density
and its total tendency.  The largest difference is a relative $L_2$ of $5.4
\times 10^{-14}$ in the current density.  Moments are reductions over the
grid, so a different summation order between the backends would explain
differences of this size; we did not verify that cause.

\begin{table}[ht]
\centering
\caption{NumPy against Rust on the parity case: relative $L_2$ difference over
all output times.}
\label{tab:parity}
\small
\begin{tabular}{lr}
\toprule
Quantity & Relative $L_2$ \\
\midrule
Distribution, every tendency, total tendency, runaway density & 0 (bit-identical) \\
Avalanche and external-source budgets of the runaway density & 0 (bit-identical) \\
Moment: current density & $5.4 \times 10^{-14}$ \\
Moment: density & $7.8 \times 10^{-16}$ \\
Moment: kinetic-energy density & $5.6 \times 10^{-16}$ \\
Runaway density: radial-transport budget & $5.2 \times 10^{-16}$ \\
Runaway density: total tendency & $2.4 \times 10^{-16}$ \\
\bottomrule
\end{tabular}
\end{table}

The declared gate of the benchmark is a relative $L_2$ of at most $2 \times
10^{-12}$ for every output.  The number of internal steps must match
exactly, and it does.

\section{Cost on one host}

Over 11 repeats on one workstation (12 logical CPUs; load averages 7.6, 14.6
and 14.6 over 1, 5 and 15 minutes; no concurrent workload declared), the
median run took 0.0126\,s with NumPy and 0.0060\,s with Rust, a ratio of
2.10.  This is an observation for this host and this small case.  It is not a
hardware-neutral speed-up, and larger grids may behave differently.

\section{Limitations}

\begin{itemize}
\item The parity case is synthetic.  Physical coefficients, larger grids and
long runs are not covered.
\item No comparison with DREAM is reported.  The repository's DREAM campaign
currently cannot run its pinned DREAM binary on the available hosts.
\item The timing comes from one loaded workstation.
\end{itemize}

\section*{Software version and date}

This preprint describes \texttt{scpn-fusion} at commit \texttt{6f22c716d099}
of 16 November 2026, on the main line after release 4.1.0.  The benchmark
report was produced at commit \texttt{98b6d4d2cfde} and tracked at
\texttt{d8a10e8fdb4a}.  The twelve sources it records are byte-identical at
the commit above: the Rust and Python grid, coefficient, operator and solver
modules, the Rust module root and Python binding, the Python diagnostics and
the benchmark.  The run used NumPy 2.5.4 and Python 3.12.3.

\section*{Data and code availability}

Code is available at \url{https://github.com/anulum/scpn-fusion-core}.  The
benchmark report cited here, with its source revision, Git blob identity and
SHA-256 digest, is in \texttt{evidence/} of this submission directory.  The
generator in that directory checks every quoted value against it.  This is a
venue-neutral preprint and has not been submitted.  No DOI has been reserved
and the manuscript-content licence has not yet been selected by the author.

\section*{Use of AI tools}

Large language model coding agents were used substantively for generating and
debugging the analysis code and for drafting and revising the manuscript text.
The agents were commercially available models, which are not named because no
vendor is involved in or sponsors this work.  They were coordinated through
SYNAPSE CHANNEL (version 0.99.52 at the time of writing;
\url{https://github.com/anulum/synapse-channel}), the open-source
agent-coordination tool of the ANULUM Agentic Coordination, Assurance and
Continuity Systems group.  The author specified the comparison, verified the
code through its tests and every reported number against the committed
report, and takes full responsibility for the content.

\bibliography{references}
\end{document}
